\documentclass[10pt]{article}

\def\Type{LessonDJ}
\def\TypeNum{Lesson 16}
\def\TypeTitle{Properties of Definite Integrals\\Notes To Instructors  }
\def\Semester{Fall 2021} %Not Used 
\def\Course{MATH 1190 \& 1210}
\def\Targets{  \bf\color{Fuchsia}I2 }

\usepackage{preambleDJ}

%\usepackage{ragged2e}


%%%%%%%%%%%% BEGIN DOCUMENT %%%%%%%%%%%%%%%
%%%%%%%%%%%% BEGIN DOCUMENT %%%%%%%%%%%%%%%
%%%%%%%%%%%% BEGIN DOCUMENT %%%%%%%%%%%%%%%
%%%%%%%%%%%% BEGIN DOCUMENT %%%%%%%%%%%%%%%
%%%%%%%%%%%% BEGIN DOCUMENT %%%%%%%%%%%%%%%




\begin{document}

\pagestyle{otherpages}
\thispagestyle{firstpage}
\vs{.65}
\lt  
\enum{
\item[I2:] Properties of Definite Integrals
\itemC{
\item Interpret and compute a definite integral in terms of signed area.
\item Compute definite integrals using their properties.
}
}

\lg Students will be able to...
\enumCN{
\item Use common shapes (rectangles, triangles, circles, etc.) to find the signed/net area between a function and the horizontal axis.
\item Reverse prior lesson goal. That is, sketch a function $f(x)$ when given  $\ds \int_a^b f(x)\,dx=c$ for various values of $[a,b]$ and $c$.
\item Use the following properties to find  the value of a definite integral.
\tasksI{2}{
\task Break Apart property
\task Reverse Limits property
\task Linearity properties
\task Area of a Line Segment property.
}
}

\feed
\itemC{
\item Many students found this lesson straight forward and appreciated the slight geometry review.
\item Some students were confused by questions on slide 6 that required using areas of common shapes to find the definite integral. This is likely due to two reasons.
\itemC{
\item  The fact that signed areas could cancel out resulting in a definite integral $= 0$, etc.
\item Some confusion about areas of common shapes (e.g. a quarter circle) and how to find radius.
}
}

 \hand
 \itemI{
 \item Page $1$ is a reminder of the meaning of signed area,  having students sketch the integrand and using areas of common shapes (triangles) to find the definite integral.
 \item Example 2/Exercise B is designed to address pre-class concerns and give more practice using common shapes to find definite integrals.
 \item Exercise C provides practice of Break Apart property
 \item Extra Practice 1 is a challenging question for most students as it requires students to reverse the process they have been engaging in and sketch the graph of $g(t)$ when given $\ds \int_a^b g(t)\,dt=c$ for several values of $[a,b]$ and $c$.
 \item Exercise D builds on work from Exercise B and provides practice for Reverse Limits property
 \item Exercise E provides practice for Linearity properties.
 \item Exercise F highlights Area of a Line Segment property.
 \item Extra Practice 2 gives more practice with combining properties to find definite integrals.
 }

\core Example 1/Exercise A, Example 2/Exercise B, Exercise C, D, E, and F.

\newpage
\important{The Definite Integral}
{
Let $a$ and $b$ be real numbers with $a<b$, and suppose $f$ is continuous on the interval $[a,b]$. The {\bf definite integral} of $f$ from $a$ to $b$ is denoted by \[\int_a^b f(x)\,dx\] and represents the {\it net area} between the graph of $f$ and the interval along the $x$-axis from $a$ to $b$. If $f$ represents the rate of change of a quantity, then $\dint_a^b f(x)\,dx$ represents the {\it net change} in the quantity as $x$ varies from $a$ to $b$.
}

%\begin{minipage}[t]{0.4\textwidth}{\bf\color{blue}Signed/Net Area}
\mpageT{.6}{.4}{
\blue{\bf Signed/Net Area for $\ds \int_a^b f(x)\;dx$}
\itemI{
\item Area of regions above the $x$-axis are {\em positive}
\item Area of regions below the $x$-axis are {\em negative}
}
\nti{
I am planning on annotating graph to include positive and negative regions on $[a,b]$.
}
}{
\begin{flushright}
\begin{tikzpicture}
  \begin{axis}[
%    axis y line = left,
%    axis x line = bottom,
 	axis line style = { line width=1pt},
    xlabel = {$x$},
    ylabel={$y=f(x)$},
    xtick       = {0.5,4.45},
    xticklabels = {$a$ ,$b$ },
    ytick       = {0},
    yticklabels = {},
    samples     = 160,
    domain      = -1:3,
    restrict y to domain=-10:100,
    xmin = -1, xmax = 5,
    ymin = -1.5, ymax = 3,
    width = 3in, height=2in
  ]
  %\addplot[green] (1,0) -- (3,0);
  %
  \addplot[name path=f, domain=-1:5,  blue, ultra thick, mark=none, ] {0.5+x*(x-2)*(x-4)/2};
  %\addplot[name path=g, domain=-1:3.5, black, thick, mark=none, ] {0.5*e^x};
  %
  \path[name path=axis] (axis cs:0,0) -- (axis cs:4.5,0);
  %\path[name path=auis] (axis cs:0,0) -- (axis cs:3,0);
  %
  \draw (0.5,0)--(0.5,1.81);
  \addplot [
        thick,
        color=green,
        fill=green, 
        fill opacity=0.1
    ]
    fill between[
        of=f and axis,
        soft clip={domain=0.5:2.25},
   ];
   \node at (1.1,.87) {\LARGE$\mathbf+$};
  \addplot [
        thick,
        color=red,
        fill=red, 
        fill opacity=0.1
    ]
    fill between[
        of=f and axis,
        soft clip={domain=2.25:3.86},
    ];
   \draw[thick] (3,-0.47)--(3.3,-0.47);
    \addplot [
        thick,
        color=green,
        fill=green, 
        fill opacity=0.1
    ]
    fill between[
        of=f and axis,
        soft clip={domain=3.86:4.45},
    ];
   \draw (4.45,0) -- (4.45,3);
   \node at (4.225,.5) {\Large$\mathbf+$};
\end{axis}
\end{tikzpicture}
\end{flushright}
}
\vs{.1}\\
{\example}\ltrm{\bf\color{Fuchsia}I2} Compute $\displaystyle\int_{-1}^2 (1-|x|)\ dx$.\\
\mpageT{.3}{.7}{
\centering
\begin{tikzpicture}
  \begin{axis}[
%    axis y line = left,
%    axis x line = bottom,
 	axis line style = { line width=1pt},
    xlabel = {$x$},
    ylabel = {$y$},
    xtick       = {-3,-2,-1,1,2,3,4,5,6},
    xticklabels = {$-3$,$-2$,$-1$,$1$,$2$,$3$,$4$,$5$,$6$},
    %ytick       = {0},
    %yticklabels = {},
    samples     = 160,
    domain      = -1:3,
    restrict y to domain=-10:100,
    xmin = -2, xmax = 2,
    ymin = -2, ymax = 2,
    width = 1.75in, height=1.75in,
    grid=both, grid style=dashed
  ]
    \addplot[name path=f, domain=-2:2,  blue, ultra thick, mark=none, ] {abs(x)};
\end{axis}
\end{tikzpicture}
\vs{-.15}
\captionof*{figure}{\small Graph of $y=|x|$}
%
\begin{tikzpicture}
  \begin{axis}[
%    axis y line = left,
%    axis x line = bottom,
 	axis line style = { line width=1pt},
    xlabel = {$x$},
    ylabel = {$y$},
    xtick       = {-3,-2,-1,1,2,3,4,5,6},
    xticklabels = {$-3$,$-2$,$-1$,$1$,$2$,$3$,$4$,$5$,$6$},
    %ytick       = {0},
    %yticklabels = {},
    samples     = 160,
    domain      = -1:3,
    restrict y to domain=-10:100,
    xmin = -2, xmax = 2,
    ymin = -2, ymax = 2,
    width = 1.75in, height=1.75in,
    grid=both, grid style=dashed
  ]
    \addplot[name path=f, domain=-2:2,  blue, ultra thick, mark=none, ] {-abs(x)};
\end{axis}
\end{tikzpicture}
\vs{-.15}
\captionof*{figure}{\small Graph of $y=-|x|$}
%
\begin{tikzpicture}
  \begin{axis}[
%    axis y line = left,
%    axis x line = bottom,
 	axis line style = { line width=1pt},
    xlabel = {$x$},
    xtick       = {-3,-2,-1,1,2,3,4,5,6},
    xticklabels = {$-3$,$-2$,$-1$,$1$,$2$,$3$,$4$,$5$,$6$},
    %ytick       = {0},
    %yticklabels = {},
    samples     = 160,
    domain      = -1:3,
    restrict y to domain=-10:100,
    xmin = -2, xmax = 2,
    ymin = -2, ymax = 2,
    width = 2in, height=2in,
    grid=both, grid style=dashed
  ]
    \addplot[name path=f, domain=-2:2,  blue, ultra thick, mark=none, ] {1-abs(x)};
      \path[name path=axis] (axis cs:-1,0) -- (axis cs:2,0);
      \addplot [
        thick,
        color=green,
        fill=green, 
        fill opacity=0.1
    ]
    fill between[
        of=f and axis,
        soft clip={domain=-1:1},
   ];
     \addplot [
        thick,
        color=red,
        fill=red, 
        fill opacity=0.1
    ]
    fill between[
        of=f and axis,
        soft clip={domain=1:2},
   ];
   \node at (-.3,.3) {\LARGE$\mathbf+$};
      \node at (.3,.3) {\LARGE$\mathbf+$};
            \node at (1.7,-.2) {\LARGE$\mathbf-$};
\end{axis}
\end{tikzpicture}
\vs{-.15}
\captionof*{figure}{\small Graph of $y=1-|x|$}
%
}{
\nti{ I will sketch and annotate graph first. Student will likely need reminders in order to sketch $f(x)=1-|x|$.
}
\sol{
After graphing $f(x)=1-|x|$ on $[-1,2]$, we  note that the region above $x$=axis on $[-1,1]$ will be positive and the region below the $x$-axis on $[1,2]$ will be negative;
\itemI{
\item The region below $f$ on $[-1,0]$ is a triangle with \\ $\ds Area_1 =\dfrac{1}{2}b\cdot h=\frac{1}{2} \cdot 1\cdot 1=\frac{1}{2}$
\item The region below $f$ on $[0,1]$ is a triangle with \\$\ds Area_2 =\dfrac{1}{2}b\cdot h=\frac{1}{2} \cdot 1\cdot 1=\frac{1}{2}$
\item The region below $f$ on $[1,2]$ is a triangle with \\$\ds Area_3 =-\dfrac{1}{2}b\cdot h=\frac{1}{2} \cdot 1\cdot 1=-\frac{1}{2}$
\item Thus $\ds \int_{-1}^2 1-|x|\; dx=Area_1+Area_2+Area_3=\frac{1}{2}+\frac{1}{2}-\frac{1}{2}=\frac{1}{2}$.
}
}
}
\newpage
{\exercise}\ltrm{\bf\color{Fuchsia}I2} Compute $\displaystyle\int_0^3(x-1)\ dx$. \\
\mpageT{.25}{.75}{
\begin{tikzpicture}
  \begin{axis}[
%    axis y line = left,
%    axis x line = bottom,
 	axis line style = { line width=1pt},
    xlabel = {$x$},
    xtick       = {-3,-2,-1,1,2,3,4,5,6},
    xticklabels = {$-3$,$-2$,$-1$,$1$,$2$,$3$,$4$,$5$,$6$},
    %ytick       = {0},
    %yticklabels = {},
    samples     = 160,
    domain      = -1:3,
    restrict y to domain=-10:100,
    xmin = -1, xmax = 3,
    ymin = -2, ymax = 2,
    width = 2in, height=2in,
    grid=both, grid style=dashed
  ]
    \addplot[name path=f, domain=-1:3,  blue, ultra thick, mark=none, ] {x-1};
      \path[name path=axis] (axis cs:0,0) -- (axis cs:3,0);
      \addplot [
        thick,
        color=green,
        fill=green, 
        fill opacity=0.1
    ]
    fill between[
        of=f and axis,
        soft clip={domain=1:3},
   ];
     \addplot [
        thick,
        color=red,
        fill=red, 
        fill opacity=0.1
    ]
    fill between[
        of=f and axis,
        soft clip={domain=0:1},
   ];
      \node at (.3,.-.2) {\LARGE$\mathbf-$};
            \node at (2,.4) {\LARGE$\mathbf+$};
\end{axis}
\end{tikzpicture}
}{
\itemC{
\mpageT{.3}{.7}[\hs{.1}]{
\item Graph $f(x)=x-1.$
}{
\item Use common shapes to find the signed (i.e. net) area between $f$ and the $x$-axis on $[0,3]$.
}
\sol{
We graph $y=f(x)=x-1$ and note:
\itemI{
\item The area below the $x$-axis on $[0,1]$ will be negative. We have a triangle with $\ds Area_1=-\frac{1}{2} b\cdot h=\frac{1}{2}\cdot 1\cdot 1=-\frac{1}{2}$.
\item The area above the $x$-axis on $[1,3]$ will be positive. We have a triangle with $\ds Area_2=\frac{1}{2} b\cdot h=\frac{1}{2}\cdot 2\cdot 2=2$.
\item $\ds \int_0^3x-1\;dx=Area_1+Area_2=-\frac{1}{2}+2=-\frac{1}{2}+\frac{4}{2}=\frac{3}{2}$
}
}
}
}


{\example}\ltrm{\bf\color{Fuchsia}I2} Use the graph of $f$ below to compute $\dint_{-1}^1 f(x)\,dx$.\\
\mpageT{.45}{.65}{
\begin{tikzpicture}
  \begin{axis}[
%    axis y line = left,
%    axis x line = bottom,
 	axis line style = { line width=1pt},
    xlabel = {$x$},
    xtick       = {-2,-1,1,2,3,4,5},
    xticklabels = {$-2$,$-1$,$1$,$2$,$3$,$4$,$5$},
        ytick       = {-2,-1,1,2},
    yticklabels = {$-2$,$-1$,$1$,$2$},
    samples     = 160,
    domain      = -1:3,
    restrict y to domain=-10:100,
    xmin = -2, xmax = 5,
    ymin = -2, ymax = 2,
        grid=both, grid style=dashed,
    width = 3.1in, height=2in
  ]
    \addplot[name path=f, domain=-2:5,  blue, ultra thick, mark=none, ] {x<0 ? 1 : (x<1 ? 1-2*x :(  x<2 ? x-2 : (x<4 ? sqrt(1-(x-3)^2)  :4-x )))};
  \path[name path=axis] (axis cs:-2,0) -- (axis cs:5,0);
     \addplot [
        thick,
        color=green,
        fill=green, 
        fill opacity=0.1
    ]
    fill between[
        of=f and axis,
        soft clip={domain=-1:0.5},
   ];
        \addplot [
        thick,
        color=red,
        fill=red, 
        fill opacity=0.1
    ]
    fill between[
        of=f and axis,
        soft clip={domain=.5:1},
   ];
           \addplot [
        thick,
        color=purple,
        fill=purple, 
        fill opacity=0.3
    ]
    fill between[
        of=f and axis,
        soft clip={domain=1:2},
   ];
   \addplot [
        thick,
        color=green,
        fill=green, 
        fill opacity=0.1
    ]
    fill between[
        of=f and axis,
        soft clip={domain=2:4},
   ];
      \addplot [
        thick,
        color=red,
        fill=red, 
        fill opacity=0.1
    ]
    fill between[
        of=f and axis,
        soft clip={domain=4:5},
   ];
	\node at (-.5,.25) {\large$\mathbf+$};
	\node at (.16,.25) {\large$\mathbf+$};
	\node at (.815,.-.1) {\large$\mathbf-$};
	\node at (1.45,.-.1) {\large$\mathbf-$};
	\node at (3,.5) {\large$\mathbf+$};
	\node at (4.75,.-.1) {\large$\mathbf-$};
\end{axis}
\end{tikzpicture}
}{
\sol{
We note that 
\itemI{
\item Area above the $x$-axis on $[-1,0.5]$ is positive. It consists of two common shapes:
\itemI{
\item On $[-1,0]$ we have a rectangle with area $Area_1=b\cdot h=1 \cdot 1=1$
\item On $[0,0.5]$ we have a triangle with area $\ds Area_2=\frac{1}{2} b \cdot h=\frac{1}{2} \cdot \frac{1}{2} \cdot 1=\frac{1}{4}$.
}
\item  Area below the $x$-axis on $[0.5,1]$ is negative. We have a triangle with $\ds Area_3=-\frac{1}{2} b \cdot h= -\frac{1}{2} \cdot \frac{1}{2} \cdot 1=-\frac{1}{4}$
\item $\ds \int_{-1}^1 f(x)\;dx = Area_1 + Area_2 +Area_3= 1 +\frac{1}{4}-\frac{1}{4}=1$
}
}
}

{\exercise}\ltrm{\bf\color{Fuchsia}I2} Use the graph of $f$ above to compute the following definite integrals:
\begin{enumerate}
\mpageT{.15}{.85}{
\item[(a)] $\dint_1^2 f(x)\,dx=$
}{ \vs{.05}
\sol{Area below the $x$-axis on $[1,2]$ is negative. We have a triangle with \\
$\ds Area=-\frac{1}{2} b \cdot h= -\frac{1}{2} \cdot 1 \cdot 1=-\frac{1}{2}$
} 
}
    ({\it Hint:} the area of a triangle of base $b$ and height $h$ is $A=\dfrac12bh$.)\\
\mpageT{.15}{.85}{
\item[(b)] $\dint_2^4 f(x)\,dx=$ 
}{ \vs{.05}
\sol{Area above the $x$-axis on $[1,2]$ is positive. We have a half-circle with \\
$\ds Area=\frac{1}{2}\pi r^2= \frac{1}{2} \pi \cdot 1^2=\frac{\pi}{2}$.  Note that the center of the half-circle is at the point $(3,0)$. We found $r$ by measuring the distance from the center to the circle at $(2,0)$ or $(4,0)$. This gives us $r=1$.
} 
}
    ({\it Hint:} the area of a circle of radius $r$ is $A=\pi r^2$.)\\
\mpageT{.15}{.85}{
\item[(c)] $\dint_1^4 f(x)\,dx=$ 
}{
\sol{ Area from a)  + Area from b) =$\ds \int_1^2 f(x)\,dx + \int_2^4 f(x)\,dx=-\frac{1}{2}+\frac{\pi}{2}$.\\\\
}
}
\blue{This illustrates the ``Break Apart" property of integrals because $$\ds \int_1^4 f(x)\,dx=\int_1^2 f(x)\,dx+\int_2^4 f(x)\,dx$$}
\end{enumerate}

\newpage
\vs{-.5}
\important{Definite Integral ``Break Apart" Property}
{
\[\dint_a^b f(x)\,dx = \dint_a^c f(x)\,dx + \dint_c^b f(x)\,dx\]
}

{\exercise}\ltrm{\bf\color{Fuchsia}I2} If $\dint_0^2 g(t)\,dt=4$, $\dint_2^5 g(t)\,dt=-3$, and $\dint_4^5 g(t)\,dt=-1$, then compute the following:
\begin{enumerate}
\mpageT{.15}{.85}{
\item[(a)] $\dint_0^5 g(t)\,dt=$
}{
\sol{ Using the Break Apart property, $\ds \int_0^5 g(t)\,dt=\int_0^2 g(t)\,dt+\int_2^5 g(t)\,dt=4+(-3)=1$
}
}
\mpageT{.15}{.85}{
\item[(b)] $\dint_0^4 g(t)\,dt=$ 
}{
\sol{ This one is trickier! We first observe
\al{
 \int_0^4 g(t)\,dt+\int_4^5 g(t)\,dt&=\int_0^5 g(t)\,dt
 \intertext{ We are told that $\ds \int_4^5 g(t)\,dt = -1$. In (a) we just computed $\int_0^5 g(t)\,dt=1$. Thus}
  \int_0^4 g(t)\,dt+ (-1)&=1\\
   \int_0^4 g(t)\,dt)&=1+1\\
     \int_0^4 g(t)\,dt&=2
 }
 More generally, we re-arranged the Break Apart property:
 $$ \int_0^4 g(t)\,dt=\int_0^5 g(t)\,dt - \int_4^5 g(t)\,dt $$
}
}

\end{enumerate}
\vs{-.1}
\begin{minipage}[t]{0.55\textwidth}
{\bf\color{blue}Extra Practice 1.}\ltrm{\bf\color{Fuchsia}I2} Sketch the graph of a continuous function $g$ satisfying all the integral equations in Exercise C. {\em Hint:} Consider using common objects with known area formulas (e.g. rectangles, triangles, etc.).
\sol{
One reasonable graph is shown on the right.  
\itemI{
\item We know that on $[0,2]$ we need an area of $4$, so we create a rectangle of height $2$ and thus \\
$\ds  \int_0^2 g(t)\,dt= b\cdot h=2\cdot 2=4$
}
}
\end{minipage}
%
\hfill
%
\begin{minipage}[t]{0.4\textwidth}
\,\\[-0.1in]
\begin{tikzpicture}
  \begin{axis}[
%    axis y line = left,
%    axis x line = bottom,
 	axis line style = { line width=1pt},
    xlabel = {$t$},
    xtick       = {-1,1,2,3,4,5,6},
    xticklabels = {$-1$,$1$,$2$,$3$,$4$,$5$,$6$},
        ytick       = {-2,-1,1,2},
    yticklabels = {$-2$,$-1$,$1$,$2$},
    ylabel ={$y=g(t)$},
        xlabel style = {anchor= west},
    ylabel style = {anchor=south west},
    samples     = 160,
    domain      = -1:3,
    restrict y to domain=-10:100,
    xmin = -1, xmax = 6,
    ymin = -2, ymax = 2,
        grid=both, grid style=dashed,
    width = 3.1in, height=2in
  ]
\addplot[name path=f, domain=0:5,  blue, ultra thick, mark=none, ] {x<2 ? 1.95 : (x<3 ? 10-4*x :(  x<4 ? -2 : -10+2*x ))};
\path[name path=axis] (axis cs:0,0) -- (axis cs:5,0);
\addplot [
        thick,
        color=green,
        fill=green, 
        fill opacity=0.1
    ]
    fill between[
        of=f and axis,
        soft clip={domain=0:2},
   ];
\addplot [
        thick,
        color=black,
        fill=black, 
        fill opacity=0.1
    ]
    fill between[
        of=f and axis,
        soft clip={domain=2:3},
   ];
   \addplot [
        thick,
        color=red,
        fill=red, 
        fill opacity=0.1
    ]
    fill between[
        of=f and axis,
        soft clip={domain=3:4},
   ];
      \addplot [
        thick,
        color=purple,
        fill=purple, 
        fill opacity=0.3
    ]
    fill between[
        of=f and axis,
        soft clip={domain=4:5},
   ];
  % \node[anchor=west] at (.05,1) {Area = $4$};
\end{axis}
\end{tikzpicture}
\end{minipage}
\blue{
\itemI{
\item On the interval $[2,3]$ we have a triangle above the $x$-axis (and thus positive area) and identical triangle below the $x$-axis (and thus negative area). Since the triangles are identical the areas cancel out and $\ds \int_2^3 g(t)\,dt=0$.
\item On $[3,4]$ we  have a rectangle below the $x$-axis, so area is negative and  $\ds \int_3^4 g(t)\,dt= -b\cdot h= -1\cdot 2=-2$.
\item  On $[4,5]$ we  have a triangle below the $x$-axis, so area is negative and  $\ds \int_4^5 g(t)\,dt= - \frac{1}{2} b\cdot h= -\frac{1}{2}\cdot  1\cdot 2=-1 $. 
\item $\ds \int_2^5 g(t)\,dt= \int_2^3 g(t)\,dt+\int_3^4 g(t)\,dt+\int_4^5 g(t)\,dt=0+(-2)+(-1)=-3$.
\item We have satisfied all integral equations: $\ds  \int_0^2 g(t)\,dt=4$, $\ds \int_2^5 g(t)\,dt=-3$ and $\ds \int_4^5 g(t)\,dt=-1$
}
}
\newpage

\important{Definite Integral ``Reverse Limits" Property}
{
\[\dint_a^b f(x)\,dx = -\dint_b^a f(x)\,dx\]
}

{\exercise}\ltrm{\bf\color{Fuchsia}I2} Use the graph of $f$ below to compute the given integrals. {\small({\it Hint:} reuse work in Exercise B!)}\\
\vs{-.3}
\mpageT{.4}{.6}[\hs{.1}]{
\begin{tikzpicture}
  \begin{axis}[
%    axis y line = left,
%    axis x line = bottom,
 	axis line style = { line width=1pt},
    xlabel = {$x$},
    xtick       = {-1,1,2,3,4,5,6},
    xticklabels = {$-1$,$1$,$2$,$3$,$4$,$5$,$6$},
        ytick       = {-2,-1,1,2},
    yticklabels = {$-2$,$-1$,$1$,$2$},
    ylabel ={$y=f(x)$},
        xlabel style = {anchor= west},
    ylabel style = {anchor=south west},
    samples     = 160,
    domain      = -1:3,
    restrict y to domain=-10:100,
    xmin = -1, xmax = 6,
    ymin = -2, ymax = 2,
        grid=both, grid style=dashed,
    width = 3.1in, height=2in
  ]
      \addplot[name path=f, domain=-2:5,  blue, ultra thick, mark=none, ] {x<0 ? 1 : (x<1 ? 1-2*x :(  x<2 ? x-2 : (x<4 ? sqrt(1-(x-3)^2)  :4-x )))};
  \path[name path=axis] (axis cs:-2,0) -- (axis cs:5,0);
\addplot [
        thick,
        color=purple,
        fill=purple, 
        fill opacity=0.3
    ]
    fill between[
        of=f and axis,
        soft clip={domain=1:2},
   ];
\addplot [
        thick,
        color=green,
        fill=green, 
        fill opacity=0.1
    ]
    fill between[
        of=f and axis,
        soft clip={domain=2:4},
   ];
	\node at (1.45,.-.1) {\large$\mathbf-$};
	\node at (3,.5) {\large$\mathbf+$};
\end{axis}
\end{tikzpicture}
}{
\mpageT{.3}{.9}{
(a) $\ds \int_2^1 f(x)\,dx=$
}{
\sol{$\ds \int_2^1 f(x)\,dx= -\int_1^2 f(x)\,dx=-\hs{-.1}\underbrace{\frac{-1}{2}}_{\text{Exercise B}}\hs{-.1}=\frac{1}{2}$
}
}
\vs{.2}
\mpageT{.3}{.9}{
(b) $\ds \int_4^2 f(x)\,dx=$
}{\sol{$\ds \int_4^2 f(x)\,dx=- \int_2^4 f(x)\,dx=-\hs{-.1} \underbrace{\frac{\pi}{2}}_{\text{Exercise B}}\hs{-.1}=-\frac{\pi}{2}$
}
}
\vs{.2}
\mpageT{.3}{.9}{
(c) $\ds \int_4^1 f(x)\,dx=$
}{
\sol{
$$ \int_4^1 f(x)\,dx=- \int_1^4 f(x)\,dx=- \underbrace{\left(-\frac{1}{2}+\frac{\pi}{2}\right)}_{\text{Exercise B}}=\frac{1}{2}-\frac{\pi}{2}$$
}
}
}


\important{Definite Integral ``Linearity" Properties}
{
\[\dint_a^b cf(x)\,dx = c\dint_a^b f(x)\,dx \quad\text{ and }\quad \dint_a^b(f(x)\pm g(x))\,dx = \dint_a^b f(x)\,dx \pm \dint_a^b g(x)\,dx\]
}

{\exercise}\ltrm{\bf\color{Fuchsia}I2} The graph of $f$ is given below. We don't know the graph of $g$, but suppose we do know $\dint_0^2 g(x)\,dx = \pi$ and $\dint_2^4 g(x)\,dx=-\pi$. Compute the following integrals.

\mpageT{.4}{.6}[\hs{.1}]{
\vspace{-0.1in}
\begin{tikzpicture}
  \begin{axis}[
	axis line style = { line width=1pt},
	xlabel = {$x$},
	xtick       = {-1,1,2,3,4,5,6},
    	xticklabels = {$-1$,$1$,$2$,$3$,$4$,$5$,$6$},
        ytick       = {-2,-1,1,2},
    	yticklabels = {$-2$,$-1$,$1$,$2$},
    	ylabel ={$y=f(x)$},
        xlabel style = {anchor= west},
    	ylabel style = {anchor=south west},
    	samples     = 160,
    	domain      = -1:3,
    	restrict y to domain=-10:100,
    	xmin = -1, xmax = 6,
    	ymin = -2, ymax = 2,
        grid=both, grid style=dashed,
    	width = 3.1in, height=2in
  ]
    \addplot[name path=f, domain=-1:6,  blue, ultra thick, mark=none, ] {x<2 ? -(4-x^2)^(1/2) :  (x<4 ? sqrt(1-(x-3)^2)  :4-x )};
  \path[name path=axis] (axis cs:-1,0) -- (axis cs:6,0);
  \addplot [
        thick,
        color=red,
        fill=red, 
        fill opacity=0.1
    ]
    fill between[
        of=f and axis,
        soft clip={domain=0:2},
   ];
     \addplot [
        thick,
        color=green,
        fill=green, 
        fill opacity=0.1
    ]
    fill between[
        of=f and axis,
        soft clip={domain=2:4},
   ];
	\node at (.9,.-.9) {\large$\mathbf-$};
	\node at (3,.5) {\large$\mathbf+$};
\end{axis}
\end{tikzpicture}
}{
(a)   $\displaystyle\int_0^2 (f(x)-g(x))\,dx$\\
\sol{
\al{ \int_0^2 (f(x)-g(x))\,dx&= \int_0^2 f(x)\,dx -\int_0^2 g(x)\,dx\\
&=-\frac{1}{4}\pi r^2 - \pi \text{ (quarter circle below $x$-axis)}\\
&=-\frac{1}{4}\pi \cdot 2^2-\pi\\
&=-\pi-\pi=-2\pi
}
}
}\vs{-.3}
\begin{enumerate}
\mpageT{.2}{.85}{
\item[(b)] $\ds \int_2^4 (f(x)-g(x))\,dx$
 }{
 \sol{$\ds \int_2^4 (f(x)-g(x))\,dx= \int_2^4 f(x)\,dx -\int_2^4 g(x)\,dx=\frac{1}{2} \pi r^2 - (-\pi)=\frac{1}{2}\pi \cdot 1^2+\pi=\frac{3\pi}{2}$
% \al{
% \int_2^4 (f(x)-g(x))\,dx&= \int_2^4 f(x)\,dx -\int_2^4 g(x)\,dx\\
% &=\frac{1}{2} \pi r^2 - (-\pi)\\
% &=\frac{1}{2}\pi \cdot 1^2+\pi\\
% &=\frac{3\pi}{2}
% }
 }
 }
 \vfill
\mpageT{.2}{.95}{
\item[(c)] $\ds \int_0^4 (f(x)-g(x))\,dx$
}{
\sol{$\ds \int_0^4 (f(x)-g(x))\,dx=\underbrace{ \int_0^2 (f(x)-g(x))\,dx+\int_2^4 (f(x)-g(x))\,dx}_{\text{Break Apart Property}}=\underbrace{-2\pi +\frac{3\pi}{2}}_{\text{See a) and b)}}=-\frac{\pi}{2}$
}
}
\vfill
\mpageT{.2}{.95}{
\item[(d)] $\ds\int_4^0 2(f(x)-g(x))\,dx$
}{
\sol{$\ds\int_4^0 2(f(x)-g(x))\,dx=2\int_4^0 (f(x)-g(x))\,dx=-2\int_0^4 (f(x)-g(x))\,dx=-2\cdot \frac{-\pi}{2}=\pi$
}
}
\end{enumerate}

\newpage

{\exercise}\ltrm{\bf\color{Fuchsia}I2} Use the graph of $f$ below to sketch a region that could represent the integral below. Find the value of the integral. \vs{-.1}\\
%\vs{-.3}
\mpageT{.4}{.6}[\hs{.1}]{
\begin{tikzpicture}
  \begin{axis}[
%    axis y line = left,
%    axis x line = bottom,
 	axis line style = { line width=1pt},
    xlabel = {$x$},
    xtick       = {-1,1,2,3,4,5,6},
    xticklabels = {$-1$,$1$,$2$,$3$,$4$,$5$,$6$},
        ytick       = {-2,-1,1,2},
    yticklabels = {$-2$,$-1$,$1$,$2$},
    ylabel ={$y=f(x)$},
        xlabel style = {anchor= west},
    ylabel style = {anchor=south west},
    samples     = 160,
    domain      = -1:3,
    restrict y to domain=-10:100,
    xmin = -1, xmax = 6,
    ymin = -2, ymax = 2,
        grid=both, grid style=dashed,
    width = 3.1in, height=2in
  ]
      \addplot[name path=f, domain=-2:5,  blue, ultra thick, mark=none, ] {x<0 ? 1 : (x<1 ? 1-2*x :(  x<2 ? x-2 : (x<4 ? sqrt(1-(x-3)^2)  :4-x )))};
  \path[name path=axis] (axis cs:-2,0) -- (axis cs:5,0);
    \draw[ultra thick, red] (3,0)--(3,1);
\end{axis}
\end{tikzpicture}
}{
\mpageT{.225}{.775}{
$\ds \int_3^3 f(x)\,dx=$
}{
\sol{$\ds \int_3^3 f(x)\,dx=0$.  The region between the $x$-axis and $f$ above $[3,3]$ is an infinitely thin line that has $0$ area. This leads to our more general ``Area of a Line Segment" property below.
}
}
}

\important{Definite Integral ``Area of a Line Segment" Property}
{
\[\dint_a^a f(x)\,dx = 0\]
}

The following problem was adapted from a problem appearing on an assessment during the Fall 2023 semester.

{\bf\color{blue}Exam Practice 2.}\ltrm{\bf\color{Fuchsia}I2} Given $\displaystyle \int_{3}^{-2} f(x)\,dx=1$, $\displaystyle \int_{0}^3 f(x)\,dx = 2$, and $\displaystyle\int_{-2}^3 \left(2f(x)-g(x)\right)\,dx=3$, evaluate the following using properties of definite integrals: 
\enumA{
\mpageT{.15}{.85}{
    \item $\ds \int_{-2}^3 f(x)\,dx$
    }{
    \sol{$\ds \int_{-2}^3 f(x)\,dx=- \int_{3}^{-2} f(x)\,dx = - (1)=-1$
    }
    }
 %   \vs{.4}
    \mpageT{.15}{.85}{
    \item $\ds \int_{-2}^0 f(x)\,dx$
 }{\vs{.05}
 \sol{\al{
  \int_{-2}^0 f(x)\,dx+ \int_{0}^3 f(x)\,dx&= \int_{-2}^3 f(x)\,dx\\
  \int_{-2}^0 f(x)\,dx+ 2&= -1\\
    \int_{-2}^0 f(x)\,dx&= -1-2\\
     \int_{-2}^0 f(x)\,dx&= -3
 }
 }
 }\vs{-.5}
 \mpageT{.15}{.85}{
    \item $\ds \int_{-2}^3 g(x)\,dx$
    }{\vs{.05}
    \sol{\al{
\int_{-2}^3 \left(2f(x)-g(x)\right)\,dx&=3\\
\int_{-2}^3 2f(x)\,dx - \int_{-2}^3 g(x)\,dx&=3\\
2\int_{-2}^3 f(x)\,dx - \int_{-2}^3 g(x)\,dx&=3\\
2(-1) - \int_{-2}^3 g(x)\,dx&=3\\
 - \int_{-2}^3 g(x)\,dx&=3+2\\
 \int_{-2}^3 g(x)\,dx&=-5
    }
    }
    }
        
}

\end{document}
